\subsection{Basic families and elementary questions}

The first examples should remain computationally simple: constants, linear functions, low-degree polynomials, reciprocal functions, absolute value, the exponential function, and the basic sine and cosine functions.

\begin{examplebox}[A quadratic function]
For
\[
f(x)=x^2-4,
\]
the domain is $\mathbb R$, the zeros are $x=\pm2$, and the function is negative between the zeros and positive outside them.
\end{examplebox}

\begin{examplebox}[A simple exponential]
For
\[
g(x)=e^x,
\]
the domain is $\mathbb R$, every value is positive, and the graph increases as $x$ increases.
\end{examplebox}

\begin{examplebox}[A periodic function]
For
\[
h(x)=\sin x,
\]
the values repeat every $2\pi$. The function oscillates between $-1$ and $1$.
\end{examplebox}

\begin{corebox}[Questions before calculus]
For an elementary function ask about its domain, values, zeros, sign, symmetry, and visible long-term behavior. Calculus will later make statements about change and limits precise.
\end{corebox}

The purpose of these examples is not to catalogue every possible function. They provide a small vocabulary rich enough to illustrate the central ideas of analysis without hiding them behind difficult algebra.
